\begin{abstract}

Modern computer vision rests on a quiet assumption that a model is handed all the visual information it needs in advance, as a complete, static, sufficiently high-resolution image. This assumption is reasonable for tasks such as classification or segmentation, but it breaks down the moment a model is placed inside an embodied agent, such as a robot or an unmanned aerial vehicle. An embodied agent perceives the world through a narrow window, with a limited field of view and resolution, must move physically to acquire each new observation, and operates under tight constraints on time, energy, and computation. Capturing and processing an entire environment at full resolution is then not merely wasteful but often infeasible. Under these conditions the central question is no longer only \emph{what we currently see}, but \emph{what we should observe} in the first place. An agent that chooses its observations wisely can match the performance of one that senses everything indiscriminately, at a fraction of the cost.

In this thesis I argue that \emph{active visual exploration} should be treated as a learning problem in its own right. I pursue this idea at two complementary scales -- the \emph{local scale} of deciding where to look within a reachable scene, and the \emph{global scale} of deciding where to move through an open world. The aim throughout is to build agents that are \emph{efficient}, understanding an environment from as few observations as possible; \emph{flexible}, exploiting the continuous capabilities of real sensors rather than artificial grids; and \emph{scalable} toward the demands of real embodied agents. This thesis comprises four publications, each of which removes a limitation left standing by the previous one.

The first contribution, \emph{Active Visual Exploration Based on Attention-Map Entropy} (IJCAI 2023), asks how an agent should decide where to look next without a separate decision network. It shows that the internal uncertainty of a transformer, read directly from the entropy of its attention maps, is itself an effective signal for selecting the most informative next glimpse, simplifying training while improving reconstruction, classification, and segmentation from partial observations. This work, however, still draws its glimpses from a fixed grid. The second contribution, \emph{Beyond Grids: Exploring Elastic Input Sampling for Vision Transformers} (WACV 2025), lifts that constraint. It formalises the notion of input \emph{elasticity}, the resilience of a model to patches of varying scale, position, and coverage, introduces a protocol to measure it, and proposes architectural and training modifications that let a vision transformer ingest patches sampled freely from an image, turning the rigid grid into a continuous sampling space. The third contribution, \emph{AdaGlimpse: Active Visual Exploration with Arbitrary Glimpse Position and Scale} (ECCV 2024), closes the gap between software and hardware. Building on an input-elastic transformer, it casts exploration as a Markov decision process and uses the Soft Actor-Critic algorithm to select glimpses of arbitrary continuous position and scale, exactly the degrees of freedom offered by optical zoom or by changing an aerial vehicle's altitude. The resulting agent discovers a coarse-to-fine strategy on its own, matching or surpassing prior methods while observing far fewer pixels.

The first three papers progressively generalise \emph{local-scale} exploration, moving from a fixed grid, through elastic patch geometry, to arbitrary continuous zoom and pan. The fourth contribution, \emph{FlySearch: Exploring how vision-language models explore} (NeurIPS 2025), turns to the complementary \emph{global-scale} problem of deciding where to move in a large, open, three-dimensional world in order to find something. It introduces a photorealistic benchmark for goal-driven object search, in which a simulated aerial vehicle must physically traverse an unbounded outdoor space to acquire each observation, and uses it to ask whether today's vision-language models can carry out this kind of exploration. It finds that state-of-the-art models fail to reliably solve even simple exploration tasks, with the gap to human performance widening sharply as tasks demand longer, consistent search strategies, and that this gap persists even after fine-tuning.

Taken together, these works trace an arc from deciding where to look within a single reachable scene to deciding where to move through an open world. They support the central claim of the thesis at the local scale, showing that exploration can be learned, can be efficient, and need not rely on sensing abstractions that diverge from current vision hardware. At the global scale, they chart the frontier rather than close it, mapping out the failure modes that define the research agenda that follows from this work.

All four publications appeared at leading venues, three at CORE A* (200 MNiSW points) conferences and one at a CORE A (140 MNiSW points) conference, and I am the first author of each. The research was supported in part by the National Science Centre Preludium grant \emph{Where to look next -- guiding active visual exploration with internal model uncertainty}, of which I was the principal investigator, and was recognised with the START scholarship of the Foundation for Polish Science. During my doctoral studies I completed a research internship at the Computer Vision group of the Max Planck Institute for Informatics under the supervision of Prof.\ Bernt Schiele. Beyond the main line of research, I contributed to work on AI-assisted medical diagnostics, which led to a European patent application, on partial multi-label learning, and on benchmarking the spatial reasoning of vision foundation models. I also served as a reviewer for NeurIPS, WACV, and KDD, and co-organised three editions of the Machine Learning Summer School in Kraków and the ELLIS Doctoral Symposium.

\vspace{1em}
\noindent\textbf{Keywords:} active visual exploration, embodied vision, active perception, vision transformers, reinforcement learning, vision-language models, object goal navigation, aerial visual search.

\end{abstract}
